\begin{ouverture}[conditionnelles]{8}{Probabilités conditionnelles et Bayes}{Probabilités et décision}{Arbres · Fréquences · Conditionnement · Prévalence}
\objectif{Traduire un énoncé avec un arbre ou un tableau.}
\objectif{Calculer une probabilité totale.}
\objectif{Inverser un conditionnement avec la formule de Bayes.}
\objectif{Interpréter sensibilité, spécificité et valeurs prédictives.}
\prerequis{Événements, intersection, réunion et probabilités de Première.}
\end{ouverture}
\essentiel

\essSection{1}{Changer l'univers de référence}
\begin{definition}
Si $P(A)>0$, la probabilité de $B$ sachant $A$ est $P_A(B)=P(A\cap B)/P(A)$.
Ainsi $P(A\cap B)=P(A)P_A(B)$.
Les événements $A$ et $B$ sont indépendants si $P(A\cap B)=P(A)P(B)$.
\end{definition}
Conditionner consiste à se restreindre aux cas où l'information est réalisée. Il ne s'agit pas simplement de multiplier deux probabilités quelconques.
Dans un arbre, les poids issus d'un même nœud totalisent $1$. Un chemin se calcule par produit ; des chemins incompatibles se réunissent par addition.
\figurelivre{c08-bayes.pdf}{Arbre d'un test fictif : $M$ désigne la présence d'une caractéristique et $+$ un résultat positif. Les fréquences à droite sont les effectifs théoriques pour $10\,000$ personnes.}
\essSection{2}{La formule des probabilités totales}
\begin{propriete}
Si $A$ et $\bar A$ ont des probabilités non nulles,
\[P(B)=P(A)P_A(B)+P(\bar A)P_{\bar A}(B).\]
Plus généralement, on additionne sur une partition de l'univers : des événements incompatibles dont la réunion est l'univers.
\end{propriete}
Un tableau d'effectifs permet souvent de vérifier les calculs de l'arbre. Les probabilités conditionnelles se lisent en divisant par le total du groupe conditionnant, pas toujours par le total général.
\essSection{3}{Inverser avec Bayes}
\begin{theoreme}[Formule de Bayes]
Pour $P(B)>0$,
\[P_B(A)=\frac{P(A)P_A(B)}{P(B)}.
\]
La probabilité $P(A)$ est dite a priori ; $P_B(A)$ est la probabilité a posteriori après observation de $B$.
\end{theoreme}
Cette formule découle directement des deux expressions de $P(A\cap B)$.
\begin{exemple}[Un test fictif]
On donne $P(M)=0{,}02$, $P_M(+)=0{,}95$ et $P_{\bar M}(+)=0{,}05$.
Alors $P(+)=0{,}02\times0{,}95+0{,}98\times0{,}05=0{,}068$.
Ainsi $P_+(M)=0{,}019/0{,}068=19/68\simeq27{,}94\,\%$.
Sur $10\,000$ personnes, on attend $190$ vrais positifs et $490$ faux positifs. Parmi les $680$ résultats positifs, $190$ concernent $M$.
\end{exemple}
\essSection{4}{Nommer les quatre groupes}
La sensibilité est $s=P_M(+)$ ; la spécificité est $c=P_{\bar M}(-)$.
La valeur prédictive positive est $P_+(M)$ ; la valeur prédictive négative est $P_-(\bar M)$.
Les vrais positifs sont $M\cap+$, les faux positifs $\bar M\cap+$ ; les faux négatifs $M\cap-$ et les vrais négatifs $\bar M\cap-$.
\begin{attention}
Une sensibilité de $95\,\%$ ne signifie pas qu'une personne ayant un test positif a $95\,\%$ de probabilité de posséder la caractéristique.
Les modèles de dépistage du manuel servent uniquement à apprendre le raisonnement probabiliste ; les données sont fictives.
\end{attention}
\essSection{5}{L'effet de la proportion initiale}
Si $p=P(M)$ et si $s,c$ sont fixés, la valeur prédictive positive vaut
\[V(p)=\frac{sp}{sp+(1-c)(1-p)}.\]
Lorsque $s>0$ et $1-c>0$, elle augmente avec $p$. Pour une caractéristique rare, des faux positifs peu fréquents peuvent être nombreux parmi tous les résultats positifs.
Deux tests successifs ne peuvent être traités comme indépendants sans hypothèse. Une indépendance conditionnelle sachant $M$, et une autre sachant $\bar M$, doivent être explicitement annoncées avant de multiplier leurs probabilités.

\exercices
\markright{Prendre confiance}
\begingroup\renewcommand{\otBesoinPlace}[1]{}\begin{minipage}{\linewidth}
\rubrique{Prendre confiance}
\begin{exercice}[1]{Lire un conditionnement}
$P(A)=0{,}4$ et $P_A(B)=0{,}7$.
\begin{questions}
\item Calculer $P(A\cap B)$.
\item Calculer $P_A(\bar B)$.
\item Peut-on déjà calculer $P(B)$ ?
\end{questions}
\end{exercice}
\begin{exercice}[1]{Compléter un arbre}
$P(A)=0{,}3$, $P_A(B)=0{,}8$ et $P_{\bar A}(B)=0{,}2$.
\begin{questions}
\item Donner les trois poids complémentaires.
\item Calculer $P(B)$.
\item Calculer $P(B\cap\bar A)$.
\end{questions}
\end{exercice}
\begin{exercice}[1]{Indépendance}
$P(A)=0{,}5$, $P(B)=0{,}4$ et $P(A\cap B)=0{,}2$.
\begin{questions}
\item Les événements sont-ils indépendants ?
\item Calculer $P_A(B)$.
\item Sont-ils incompatibles ?
\end{questions}
\end{exercice}
\begin{exercice}[1]{Deux urnes}
On choisit une urne $U$ avec probabilité $1/3$, sinon $V$. $U$ contient $3$ boules rouges sur $4$ ; $V$ en contient $1$ sur $4$.
\begin{questions}
\item Calculer $P(R)$.
\item Calculer $P_R(U)$.
\item Comparer $P_R(U)$ à $P(U)$.
\end{questions}
\end{exercice}
\end{minipage}\endgroup
\clearpage
\markright{Appliquer et relier}
\begingroup\renewcommand{\otBesoinPlace}[1]{}\begin{minipage}{\linewidth}
\rubrique{Appliquer et relier}
\begin{exercice}[2]{Production et défauts}
Une machine $A$ produit $60\,\%$ des pièces, avec $2\,\%$ de défauts ; $B$ produit le reste, avec $5\,\%$ de défauts.
\begin{questions}
\item Calculer le taux global de défaut.
\item Une pièce est défectueuse : calculer la probabilité qu’elle vienne de $B$.
\item Pourquoi la moyenne simple de $2\,\%$ et $5\,\%$ serait-elle incorrecte ?
\end{questions}
\end{exercice}
\begin{exercice}[2]{Un tableau de test}
Sur $1\,000$ personnes fictives, $50$ possèdent $M$. Parmi elles, $45$ ont un résultat positif. Parmi les autres, $19$ ont un résultat positif.
\begin{questions}
\item Donner sensibilité et spécificité.
\item Calculer la valeur prédictive positive.
\item Calculer la valeur prédictive négative.
\end{questions}
\end{exercice}
\begin{exercice}[2]{Trois origines}
Trois ateliers ont les parts $0{,}2$, $0{,}3$, $0{,}5$ et les taux de défaut $0{,}01$, $0{,}02$, $0{,}04$.
\begin{questions}
\item Calculer la probabilité d’un défaut.
\item Calculer la part du troisième atelier parmi les pièces défectueuses.
\item Vérifier que les trois probabilités a posteriori totalisent $1$.
\end{questions}
\end{exercice}
\begin{exercice}[2]{Retrouver une proportion}
$P_A(B)=0{,}8$, $P_{\bar A}(B)=0{,}2$ et $P(B)=0{,}35$.
\begin{questions}
\item Poser $p=P(A)$ et écrire l’équation.
\item Déterminer $p$.
\item Calculer $P_B(A)$.
\end{questions}
\end{exercice}
\end{minipage}\endgroup
\clearpage
\markright{Chercher et argumenter}
\begingroup\renewcommand{\otBesoinPlace}[1]{}\begin{minipage}{\linewidth}
\rubrique{Chercher et argumenter}
\begin{exercice}[3]{Prévalence et valeur prédictive}
Un test fictif a $s=0{,}9$ et $c=0{,}95$.
\begin{questions}
\item Exprimer $V(p)$.
\item Déterminer les $p$ pour lesquels $V(p)>1/2$.
\item Interpréter sans confondre les deux conditionnements.
\end{questions}
\end{exercice}
\begin{exercice}[3]{Deux résultats positifs}
$P(M)=0{,}1$, $P_M(+)=0{,}9$, $P_{\bar M}(+)=0{,}1$. On suppose deux répétitions indépendantes conditionnellement à chacun des deux groupes.
\begin{questions}
\item Calculer la probabilité de deux positifs sachant chaque groupe.
\item Calculer la probabilité de $M$ sachant deux positifs.
\item Que faudrait-il revoir si les erreurs des répétitions étaient liées ?
\end{questions}
\end{exercice}
\begin{exercice}[3]{L’erreur de l’inversion}
Un élève affirme : « $P_A(B)=0{,}8$, donc $P_B(A)=0{,}8$. »
\begin{questions}
\item Construire un contre-exemple avec $P(A)=0{,}1$ et $P_{\bar A}(B)=0{,}2$.
\item Calculer les deux probabilités.
\item Formuler la règle correcte.
\end{questions}
\end{exercice}
\begin{exercice}[3]{Une preuve de croissance}
$V(p)=sp/(sp+f(1-p))$ avec $s>0,f>0$ et $0\leq p\leq1$.
\begin{questions}
\item Calculer la dérivée.
\item En déduire le sens de variation.
\item Donner les valeurs aux bornes et une conclusion.
\end{questions}
\end{exercice}
\end{minipage}\endgroup
